% --- Archon physics formalization source begin ---
% archon:physics
% archon:covers PhyXMiniProblems/problem_phyx_mini_0052.lean
% archon:source-report reports/phyx_mini/problem_phyx_mini_0052.source.json

\chapter{Physics problem phyx\_mini\_0052}
\label{ch:PhyXMiniProblems_problem_phyx_mini_0052}

\paragraph{Problem source.}
\#\# Physical scenario

To see a flower better, a naturalist holds a 6.0-cm-focal-length magnifying glass 4.0 cm from the flower.

\#\# Question

What is the magnification?

\#\# Answer choices

- A: A: 6.4

- B: B: 3.0

- C: C: 8.9

- D: D: 1.24

\#\# Figure caption (auxiliary; use the image as primary evidence)

The image depicts a ray diagram illustrating the formation of an image by a converging lens. Key elements and features include:

- **Objects and Image**: 
  - An upright green arrow near the center labeled "Object," located at a distance of 4 cm from the lens.
  - A larger upright orange arrow labeled "Image," positioned at a distance of 12 cm from the lens on the opposite side.

- **Lens**:
  - A convex (converging) lens is represented in the center, with its edges curved outward.

- **Focal Points**:
  - Two focal points labeled "Focal point" are marked on each side of the lens, each 4 cm from the lens.

- **Rays**:
  - Solid and dashed purple lines represent light rays. 
  - The rays diverge from the top of the object, pass through the lens, and then converge to form the image on the other side.
  - Dotted purple lines trace the rays backward to locate the image.

- **Dimensional Labels**:
  - The object distance (s) is marked as 4 cm, and the image distance (s') as 12 cm.
  - The focal length (f) is indicated as 4 cm.

- **Instructional Text**:
  - Blue text on the right side says, "Trace these rays back to the image location."

- **Axes**:
  - A horizontal line serves as the principal axis, with distances marked along it.

Overall, the diagram visually explains how light rays passing through a convex lens create a magnified image at a specific location.

\#\# Dataset metadata

Geometrical Optics; Physical Model Grounding Reasoning, Multi-Formula Reasoning

\paragraph{Recorded answer/context.}
B: B: 3.0

\paragraph{Figure/image path.}
/root/proposal\_for\_physic/science-mango/phyx\_mini\_run/phyx\_data/test\_image/52.png

\paragraph{Formalization target.}
create a compiling Lean file with sorry bodies at `PhyXMiniProblems/problem\_phyx\_mini\_0052.lean`. The Lean declarations must preserve the physical quantities, dimensions or dimensional roles, figure labels, governing-law hypotheses, and final relation expressed by this problem.
Use Mathlib/Physlib names found through LeanExplore where available. If a physics API is missing, introduce faithful local abstractions rather than scalar placeholder aliases.

\begin{theorem}[Physics formalization target]
\label{thm:physics:phyx_mini_0052:target}
\lean{PhyXMiniProblems.ProblemPhyXMini0052.problem_phyx_mini_0052}
\uses{def:physics:phyx-mini-0052:phyxminiproblems-problemphyxmini0052-magnifyingglassdiagram, def:physics:phyx-mini-0052:phyxminiproblems-problemphyxmini0052-hasstatedproblemdata, def:physics:phyx-mini-0052:phyxminiproblems-problemphyxmini0052-haspositiveproblemlengths, def:physics:phyx-mini-0052:phyxminiproblems-problemphyxmini0052-matchesraydiagram, def:physics:phyx-mini-0052:phyxminiproblems-problemphyxmini0052-usescartesiansignconvention, def:physics:phyx-mini-0052:phyxminiproblems-problemphyxmini0052-obeysthinlensequation, def:physics:phyx-mini-0052:phyxminiproblems-problemphyxmini0052-obeystransversemagnificationlaw, def:physics:phyx-mini-0052:phyxminiproblems-problemphyxmini0052-iscorrectanswer}
The assigned autoformalize agent should translate this physics problem into Lean declarations in the covered file, with theorem and lemma proof bodies written as `by sorry`.
\end{theorem}
\begin{proof}
This is an autoformalization task, not a proof task. Produce statements that can later be proved by the physics prover without weakening the physical model.
\end{proof}

% --- Archon declaration topology begin ---
\section{Declaration topology}
\begin{definition}[Optical Length]
  \label{def:physics:phyx-mini-0052:phyxminiproblems-problemphyxmini0052-opticallength}
  \lean{PhyXMiniProblems.ProblemPhyXMini0052.OpticalLength}
  A signed physical length, independent of the unit used to read it
\end{definition}
\begin{proof}
  This is the declared model definition of Optical Length.
\end{proof}
\begin{definition}[centimeter Unit Choices]
  \label{def:physics:phyx-mini-0052:phyxminiproblems-problemphyxmini0052-centimeterunitchoices}
  \lean{PhyXMiniProblems.ProblemPhyXMini0052.centimeterUnitChoices}
  Unit choices in which the length coordinate is read in centimeters
\end{definition}
\begin{proof}
  This is the declared model definition of centimeter Unit Choices.
\end{proof}
\begin{definition}[centimeters Value]
  \label{def:physics:phyx-mini-0052:phyxminiproblems-problemphyxmini0052-centimetersvalue}
  \lean{PhyXMiniProblems.ProblemPhyXMini0052.centimetersValue}
  \uses{def:physics:phyx-mini-0052:phyxminiproblems-problemphyxmini0052-opticallength, def:physics:phyx-mini-0052:phyxminiproblems-problemphyxmini0052-centimeterunitchoices}
  The signed real-valued centimeter readout of a physical length
\end{definition}
\begin{proof}
  This is the declared model definition of centimeters Value.
\end{proof}
\begin{definition}[Thin Lens Kind]
  \label{def:physics:phyx-mini-0052:phyxminiproblems-problemphyxmini0052-thinlenskind}
  \lean{PhyXMiniProblems.ProblemPhyXMini0052.ThinLensKind}
  The lens type identified by the convex profile in the ray diagram
\end{definition}
\begin{proof}
  This is the declared model definition of Thin Lens Kind.
\end{proof}
\begin{definition}[Image Nature]
  \label{def:physics:phyx-mini-0052:phyxminiproblems-problemphyxmini0052-imagenature}
  \lean{PhyXMiniProblems.ProblemPhyXMini0052.ImageNature}
  Whether the rays themselves meet at the image or only their backward extensions do
\end{definition}
\begin{proof}
  This is the declared model definition of Image Nature.
\end{proof}
\begin{definition}[Image Orientation]
  \label{def:physics:phyx-mini-0052:phyxminiproblems-problemphyxmini0052-imageorientation}
  \lean{PhyXMiniProblems.ProblemPhyXMini0052.ImageOrientation}
  Orientation of the image arrow relative to the flower/object arrow
\end{definition}
\begin{proof}
  This is the declared model definition of Image Orientation.
\end{proof}
\begin{definition}[Magnifying Glass Diagram]
  \label{def:physics:phyx-mini-0052:phyxminiproblems-problemphyxmini0052-magnifyingglassdiagram}
  \lean{PhyXMiniProblems.ProblemPhyXMini0052.MagnifyingGlassDiagram}
  \uses{def:physics:phyx-mini-0052:phyxminiproblems-problemphyxmini0052-opticallength, def:physics:phyx-mini-0052:phyxminiproblems-problemphyxmini0052-thinlenskind, def:physics:phyx-mini-0052:phyxminiproblems-problemphyxmini0052-imagenature, def:physics:phyx-mini-0052:phyxminiproblems-problemphyxmini0052-imageorientation}
  The physical quantities and labels in the magnifying-glass ray diagram. objectDistanceFromLens is positive for the flower on the incoming side, whereas signedImageDistanceFromLens is negative for the virtual image in this problem. The two focal-point positions and the object/image positions retain the geometry shown on the principal axis
\end{definition}
\begin{proof}
  This is the declared model definition of Magnifying Glass Diagram.
\end{proof}
\begin{definition}[Has Stated Problem Data]
  \label{def:physics:phyx-mini-0052:phyxminiproblems-problemphyxmini0052-hasstatedproblemdata}
  \lean{PhyXMiniProblems.ProblemPhyXMini0052.HasStatedProblemData}
  \uses{def:physics:phyx-mini-0052:phyxminiproblems-problemphyxmini0052-centimetersvalue, def:physics:phyx-mini-0052:phyxminiproblems-problemphyxmini0052-magnifyingglassdiagram}
  Problem data: a converging magnifying glass of focal length 6.0 cm is held 4.0 cm from the flower
\end{definition}
\begin{proof}
  This is the declared model definition of Has Stated Problem Data.
\end{proof}
\begin{definition}[Has Positive Problem Lengths]
  \label{def:physics:phyx-mini-0052:phyxminiproblems-problemphyxmini0052-haspositiveproblemlengths}
  \lean{PhyXMiniProblems.ProblemPhyXMini0052.HasPositiveProblemLengths}
  \uses{def:physics:phyx-mini-0052:phyxminiproblems-problemphyxmini0052-centimetersvalue, def:physics:phyx-mini-0052:phyxminiproblems-problemphyxmini0052-magnifyingglassdiagram}
  The stated focal length and flower-to-lens distance are physical positive lengths
\end{definition}
\begin{proof}
  This is the declared model definition of Has Positive Problem Lengths.
\end{proof}
\begin{definition}[Matches Ray Diagram]
  \label{def:physics:phyx-mini-0052:phyxminiproblems-problemphyxmini0052-matchesraydiagram}
  \lean{PhyXMiniProblems.ProblemPhyXMini0052.MatchesRayDiagram}
  \uses{def:physics:phyx-mini-0052:phyxminiproblems-problemphyxmini0052-centimetersvalue, def:physics:phyx-mini-0052:phyxminiproblems-problemphyxmini0052-magnifyingglassdiagram}
  Geometry and qualitative labels read from the ray diagram. The numerical image location is deliberately absent: the dashed rays identify the image as virtual and upright, but its -12 cm signed distance is derived from the lens law below
\end{definition}
\begin{proof}
  This is the declared model definition of Matches Ray Diagram.
\end{proof}
\begin{definition}[Uses Cartesian Sign Convention]
  \label{def:physics:phyx-mini-0052:phyxminiproblems-problemphyxmini0052-usescartesiansignconvention}
  \lean{PhyXMiniProblems.ProblemPhyXMini0052.UsesCartesianSignConvention}
  \uses{def:physics:phyx-mini-0052:phyxminiproblems-problemphyxmini0052-centimetersvalue, def:physics:phyx-mini-0052:phyxminiproblems-problemphyxmini0052-magnifyingglassdiagram}
  The Cartesian sign convention used in the figure: a virtual image has negative signed image distance, and an upright image has positive magnification
\end{definition}
\begin{proof}
  This is the declared model definition of Uses Cartesian Sign Convention.
\end{proof}
\begin{definition}[Obeys Thin Lens Equation]
  \label{def:physics:phyx-mini-0052:phyxminiproblems-problemphyxmini0052-obeysthinlensequation}
  \lean{PhyXMiniProblems.ProblemPhyXMini0052.ObeysThinLensEquation}
  \uses{def:physics:phyx-mini-0052:phyxminiproblems-problemphyxmini0052-magnifyingglassdiagram}
  The paraxial thin-lens equation 1/f = 1/s + 1/s', written in the division-free, dimensionally homogeneous form f * (s + s') = s * s'. Requiring it for every unit choice makes the governing law independent of the centimeter readout used for the numerical data
\end{definition}
\begin{proof}
  This is the declared model definition of Obeys Thin Lens Equation.
\end{proof}
\begin{definition}[Obeys Transverse Magnification Law]
  \label{def:physics:phyx-mini-0052:phyxminiproblems-problemphyxmini0052-obeystransversemagnificationlaw}
  \lean{PhyXMiniProblems.ProblemPhyXMini0052.ObeysTransverseMagnificationLaw}
  \uses{def:physics:phyx-mini-0052:phyxminiproblems-problemphyxmini0052-magnifyingglassdiagram}
  The signed transverse-magnification law m = -s'/s
\end{definition}
\begin{proof}
  This is the declared model definition of Obeys Transverse Magnification Law.
\end{proof}
\begin{definition}[Has Displayed Image Location]
  \label{def:physics:phyx-mini-0052:phyxminiproblems-problemphyxmini0052-hasdisplayedimagelocation}
  \lean{PhyXMiniProblems.ProblemPhyXMini0052.HasDisplayedImageLocation}
  \uses{def:physics:phyx-mini-0052:phyxminiproblems-problemphyxmini0052-centimetersvalue, def:physics:phyx-mini-0052:phyxminiproblems-problemphyxmini0052-magnifyingglassdiagram}
  The numerical image-location annotation visible in the auxiliary ray diagram. This is a consistency property to be derived, not a premise of the target
\end{definition}
\begin{proof}
  This is the declared model definition of Has Displayed Image Location.
\end{proof}
\begin{definition}[Answer Choice]
  \label{def:physics:phyx-mini-0052:phyxminiproblems-problemphyxmini0052-answerchoice}
  \lean{PhyXMiniProblems.ProblemPhyXMini0052.AnswerChoice}
  Labels of the four displayed multiple-choice answers
\end{definition}
\begin{proof}
  This is the declared model definition of Answer Choice.
\end{proof}
\begin{definition}[answer Magnification]
  \label{def:physics:phyx-mini-0052:phyxminiproblems-problemphyxmini0052-answermagnification}
  \lean{PhyXMiniProblems.ProblemPhyXMini0052.answerMagnification}
  \uses{def:physics:phyx-mini-0052:phyxminiproblems-problemphyxmini0052-answerchoice}
  The dimensionless magnification printed beside each answer choice
\end{definition}
\begin{proof}
  This is the declared model definition of answer Magnification.
\end{proof}
\begin{definition}[Is Correct Answer]
  \label{def:physics:phyx-mini-0052:phyxminiproblems-problemphyxmini0052-iscorrectanswer}
  \lean{PhyXMiniProblems.ProblemPhyXMini0052.IsCorrectAnswer}
  \uses{def:physics:phyx-mini-0052:phyxminiproblems-problemphyxmini0052-magnifyingglassdiagram, def:physics:phyx-mini-0052:phyxminiproblems-problemphyxmini0052-answerchoice, def:physics:phyx-mini-0052:phyxminiproblems-problemphyxmini0052-answermagnification}
  A choice is correct precisely when its printed value equals the physical magnification
\end{definition}
\begin{proof}
  This is the declared model definition of Is Correct Answer.
\end{proof}
\begin{lemma}[signed Image Distance eq neg twelve]
  \label{lem:physics:phyx-mini-0052:phyxminiproblems-problemphyxmini0052-signedimagedistance-eq-neg-twelve}
  \lean{PhyXMiniProblems.ProblemPhyXMini0052.signedImageDistance_eq_neg_twelve}
  \uses{def:physics:phyx-mini-0052:phyxminiproblems-problemphyxmini0052-centimetersvalue, def:physics:phyx-mini-0052:phyxminiproblems-problemphyxmini0052-magnifyingglassdiagram, def:physics:phyx-mini-0052:phyxminiproblems-problemphyxmini0052-hasstatedproblemdata, def:physics:phyx-mini-0052:phyxminiproblems-problemphyxmini0052-obeysthinlensequation}
  The thin-lens law and the two stated distances determine the virtual-image distance s' = -12 cm; this also matches the image label in the ray diagram
\end{lemma}
\begin{proof}
  The conclusion follows from the governing relations and figure readouts named in the statement of signed Image Distance eq neg twelve.
\end{proof}
\begin{lemma}[displayed Image Location follows]
  \label{lem:physics:phyx-mini-0052:phyxminiproblems-problemphyxmini0052-displayedimagelocation-follows}
  \lean{PhyXMiniProblems.ProblemPhyXMini0052.displayedImageLocation_follows}
  \uses{def:physics:phyx-mini-0052:phyxminiproblems-problemphyxmini0052-magnifyingglassdiagram, def:physics:phyx-mini-0052:phyxminiproblems-problemphyxmini0052-hasstatedproblemdata, def:physics:phyx-mini-0052:phyxminiproblems-problemphyxmini0052-matchesraydiagram, def:physics:phyx-mini-0052:phyxminiproblems-problemphyxmini0052-obeysthinlensequation, def:physics:phyx-mini-0052:phyxminiproblems-problemphyxmini0052-hasdisplayedimagelocation}
  The problem data, thin-lens equation, and labeled axis geometry reproduce the auxiliary diagram's image position 12 cm to the left of the lens
\end{lemma}
\begin{proof}
  The conclusion follows from the governing relations and figure readouts named in the statement of displayed Image Location follows.
\end{proof}
% --- Archon declaration topology end ---
% --- Archon physics formalization source end ---
